\diarytitle{0024}{積分因子による完全化と一般解（μ(x) の場合）}

$M(x,y)\,dx + N(x,y)\,dy = 0$ が完全でない（$M_y \neq N_x$）場合でも、適当な関数 $\mu$（積分因子）を掛けた $\mu M\,dx + \mu N\,dy = 0$ が完全微分方程式になることがある。もし $\dfrac{M_y - N_x}{N}$ が $x$ のみの関数ならば、$\mu = \mu(x)$ ととれて
\[
\mu(x) = \exp\left(\int \frac{M_y - N_x}{N}\,dx\right)
\]
と求まる。積分因子を掛けた後は、通常の完全微分方程式と同様に $F(x,y) = C$ を求めればよい。

$M = 3xy + y^{2}$、$N = x^{2} + xy$ とおくと
\[
\frac{\partial M}{\partial y} = 3x + 2y, \qquad \frac{\partial N}{\partial x} = 2x + y
\]
一致しないので完全ではない。
\[
\frac{M_y - N_x}{N} = \frac{(3x+2y)-(2x+y)}{x^{2}+xy} = \frac{x+y}{x(x+y)} = \frac{1}{x}
\]
は $x$ のみの関数だから、$x$ のみに依存する積分因子が存在し
\[
\mu(x) = \exp\left(\int \frac{1}{x}\,dx\right) = e^{\ln x} = x
\]
両辺に $\mu(x) = x$ を掛けると
\[
(3x^{2}y + xy^{2})\,dx + (x^{3} + x^{2}y)\,dy = 0
\]
新しい $M^{*} = 3x^{2}y + xy^{2}$、$N^{*} = x^{3} + x^{2}y$ について
\[
\frac{\partial M^{*}}{\partial y} = 3x^{2} + 2xy, \qquad \frac{\partial N^{*}}{\partial x} = 3x^{2} + 2xy
\]
一致し完全である。$F_x = M^{*}$ より
\[
F = \int (3x^{2}y + xy^{2})\,dx = x^{3}y + \frac{x^{2}y^{2}}{2} + g(y)
\]
$F_y = x^{3} + x^{2}y + g'(y) = N^{*} = x^{3} + x^{2}y$ より $g'(y) = 0$。よって一般解は
\[
\boxed{\; x^{3}y + \frac{x^{2}y^{2}}{2} = C \;}
\qquad (C \text{ は任意定数})
\]
（検算: $F_x = 3x^{2}y + xy^{2} = M^{*}$、$F_y = x^{3} + x^{2}y = N^{*}$ で一致する。）
